\documentclass{amsart}
\usepackage{amssymb,mathtools,array,longtable}
\usepackage[hidelinks]{hyperref}
\hypersetup{pdftitle={Irrationality of the dilogarithm at 1/4},pdfauthor={Ruiyi Zhang}}

\makeatletter
\def\LT@makecaption#1#2#3{\LT@mcol\LT@cols c{\hbox to\z@{\hss
  \parbox[t]\textwidth{\normalsize\@tempcnta=64 \@makecaption{#2}{#3}}\hss}}}
\makeatother

\newtheorem{theorem}{Theorem}
\newtheorem{lemma}{Lemma}[section]
\newtheorem{proposition}[lemma]{Proposition}
\numberwithin{equation}{section}
\numberwithin{table}{section}

\newcommand{\Li}{\operatorname{Li}}
\newcommand{\Res}{\operatorname*{Res}}
\newcommand{\Z}{\mathbb{Z}}
\newcommand{\fl}[1]{\lfloor #1\rfloor}
\newcommand{\br}[1]{[\,#1\,]}

\begin{document}

\title{Irrationality of the dilogarithm at 1/4}
\author{Ruiyi Zhang}
\address{Independent Researcher, Columbia, Maryland, USA}
\email{ryzfsu@gmail.com}
\subjclass[2020]{Primary 11J72; Secondary 33B30}
\keywords{dilogarithm, irrationality, permutation group method}

\begin{abstract}
We prove that $\operatorname{Li}_2(1/4)$ is irrational. Irrationality of
$\operatorname{Li}_2(1/s)$ was known for integers $s\le-5$ and $s\ge6$ by
work of Hata and of Rhin and Viola, and Prellberg recently proved it for
$s=-4,5,-3$. The linear forms in
$1$ and $\operatorname{Li}_2(1/4)$ come from one double integral of Viola and
Zudilin at $z=4$. We expand it into Rhin--Viola integrals in three ways and
apply their factorial transformations to each term. This gives a common factor
of the coefficients of the $N$th form. The third expansion adds the primes
between $265N$ and $322N$ to this factor, which is too small without them. We
bound the integral by the weighted arithmetic-geometric mean inequality and by
concavity on circular arcs.
\end{abstract}

\maketitle

\section{Introduction}

For $|w|\le1$ let $\Li_2(w)=\sum_{n\ge1}w^n/n^2$. Hata \cite{Ha} proved that
$\Li_2(1/s)$ is irrational for every integer $s\le-5$ and every $s\ge7$. Rhin
and Viola adapted their permutation group method \cite{RV96} to the dilogarithm
and settled $s=6$ \cite{RV05}. They later treated negative arguments
\cite{RV19}. Viola and Zudilin \cite{VZ} introduced a six-parameter family of
double integrals. In 2024 Calegari, Dimitrov and Tang
\cite[Remark~2.8.2 and Section~15.2]{CDT} listed the integers still open:
$s\in\{-4,-3,-2,2,3,4,5\}$. Prellberg's 2026 preprint \cite{Pr} treats
$s=-4,5,-3$. We prove the case $s=4$.

\begin{theorem}\label{thm:main}
The number $\Li_2(1/4)$ is irrational.
\end{theorem}

By Landen's identity \cite[(25.12.3)]{DLMF},
$\Li_2(1/4)+\Li_2(-1/3)=-\frac12\log^2(3/4)$, so Theorem~\ref{thm:main} does
not follow directly from the case $s=-3$.

Put $\theta=\Li_2(1/4)$, $d_m=\operatorname{lcm}(1,\dots,m)$, $d_0=1$, and
for $N\ge1$
\begin{equation}\label{eq:GD}
\begin{aligned}
G_N&=J_4(240N,240N,190N,157N,215N,165N),\\
D_N&=d_{323N}\,d_{322N}\,4^{132N}\,3^{107N},
\end{aligned}
\end{equation}
with $J_4$ as in Section~\ref{sec:int}. By \cite[Lemma~2.1]{VZ},
$D_NG_N\in\Z+\Z\theta$. Sections~\ref{sec:div} and~\ref{sec:dens} give
$\Delta_N\mid D_N$ with
\begin{gather*}
(D_N/\Delta_N)\,G_N\in\Z+\Z\theta,\\
\lim_{N\to\infty}\frac{\log\Delta_N}{40N}=\delta>4.7670,\qquad
\lim_{N\to\infty}\frac{\log D_N}{40N}<23.63856,
\end{gather*}
and Sections~\ref{sec:an} and~\ref{sec:proof} give
\[
G_N>0\ \text{for large }N,\qquad
\limsup_{N\to\infty}\frac{\log G_N}{40N}\le-18.87333 .
\]
Since $23.63856-4.7670-18.87333<0$, Theorem~\ref{thm:main} follows.

The arithmetic follows \cite[Section~3.2]{Pr}. Identity \eqref{eq:E1} for
$G_N$ and for its image under (F2) gives the analogues of
\cite[(3.11)--(3.12)]{Pr}. The factorial transformations of \cite{RV05} act on
each term. For $265N<p<322N$ these expansions give only $M=0$ in
Lemma~\ref{lem:gains}. These primes account for $57/40$ of $\delta$, and
the rest of $\delta$ is below $3.35<23.63856-18.87333$. Expansion
\eqref{eq:E3} gives one power of each such $p$. The bounds of
Lemma~\ref{lem:gains} depend only on $N/p$, so we need not exclude the
exceptional primes of \cite[(3.6)--(3.7)]{Pr}. One linear form suffices,
since $G_N>0$ for large $N$.

Two Python scripts recheck only part of the proof (Appendix~\ref{app:ver}).
The proof is complete without them. The use of AI tools is described at the
end of the paper.

\section{Integrals and known results}\label{sec:int}

Let $z>1$. For integers $h,j,k,l,m,q\ge0$ with
$j+k-m,\ j+q-m,\ h+m-k,\ h+q-k\ge0$ put
\[
R_z(x,y)=z^{k-l-q}\,
\frac{x^j(1-x)^h\,y^k(1-y)^l\,(1-y+zy)^{j+q-m}}
{\bigl(x(1-y)+zy\bigr)^{j+k-m+1}},
\]
and, with small positively oriented circles,
\begin{align*}
J_z^{(0)}&=\int_0^1\!\!\int_0^1R_z\,dy\,dx,&
J_z^{(1)}&=\int_0^1\frac1{2\pi i}\oint_{|y-\frac{x}{x-z}|=\varepsilon}R_z\,dy\,dx,\\
J_z^{(2)}&=\frac1{2\pi i}\oint_{|x-z|=\varepsilon'}\frac1{2\pi i}
\oint_{|y-\frac{x}{x-z}|=\varepsilon}R_z\,dy\,dx,&
J_z&=J_z^{(0)}-(\log z)J_z^{(1)}.
\end{align*}
For $t=(h,j,k,l,m)\in\Z_{\ge0}^5$ define $I_z^{(0)}(t)$, $I_z^{(1)}(t)$,
$I_z^{(2)}(t)$ and $I_z(t)$ in the same way from
\[
z^{-l-m}\,\frac{x^j(1-x)^h\,y^k(1-y)^l}{\bigl(x(1-y)+zy\bigr)^{j+k-m+1}}
\qquad\text{\cite[(2.1)--(2.4)]{RV05}}.
\]
For $\mathcal J=J_z$ or $\mathcal J=I_z(t)$, with components
$\mathcal J^{(0)},\mathcal J^{(1)},\mathcal J^{(2)}$, put
\[
c(\mathcal J)=\bigl(c_1(\mathcal J),c_2(\mathcal J)\bigr)
=\bigl(\mathcal J+\Li_2(1/z)\,\mathcal J^{(2)},\ \mathcal J^{(2)}\bigr).
\]
For $t=(h,j,k,l,m)$ let
\begin{gather*}
u(t)=(l+m-j,\ h+m-k,\ h+j-l,\ j+k-m),\qquad T(t)=h!\,j!\,k!\,l!\,m!,\\
H(t)=\max_iu_i,\qquad K(t)=\max\{u_1,\min(u_2,u_3),u_4\},
\end{gather*}
and call $t$ \emph{admissible} if $u(t)\in\Z_{\ge0}^4$. Put
\[
\rho(t)=(h,m,l,k,j),\qquad \sigma(t)=(h+m-k,\ j+k-m,\ m,\ l,\ k),
\]
and let $O(t)$ be the orbit of $t$ under the group generated by $\rho$ and
$\sigma$. Both maps permute the multiset
$\{h,j,k,l,m,u_1,u_2,u_3,u_4\}$ \cite[(3.2), (3.15)]{RV05}, so every
$s\in O(t)$ is admissible when $t$ is.

\begin{enumerate}
\item[(F1)] \cite[Lemma 2.1]{VZ} Let
$H_1=\max\{l+q,\,h+m-k,\,h+j-l,\,j+k-m\}$,
$H_2=\max\{l+q,\,\min(h+m-k,h+j-l),\,j+k-m\}$,
$\alpha=\max\{0,l+q-k,l+q-m\}$ and $\beta=\max\{0,k+l-h\}$. There are
$P,Q\in\Z[z]$ with
\[
d_{H_1}d_{H_2}\,z^{\alpha}(1-z)^{\beta}\,c(J_z)=(P(z),Q(z)) .
\]
\item[(F2)] \cite[Lemma 3.1]{VZ} $c(J_z(h,j,k,l,m,q))=c(J_z(j,h,m,l,k,q))$.
\item[(F3)] \cite[Theorem 2.1 and (2.9)]{RV05} Let
$\alpha(t)=\max\{j+k,k+l,l+m\}$ and $\beta(t)=\max\{0,k+l-h\}$. There are
$P_t,Q_t\in\Z[z]$ with
\[
d_{H(t)}d_{K(t)}\,z^{\alpha(t)}(z-1)^{\beta(t)}\,c(I_z(t))=(P_t(z),Q_t(z)) .
\]
\item[(F4)] \cite[(3.10)--(3.16)]{RV05} If $t$ is admissible and $s\in O(t)$,
then $c(I_z(t))/T(t)=c(I_z(s))/T(s)$.
\end{enumerate}

In the sources, (F1) and (F3) read $\varphi\mathcal J=P-Q\Li_2(1/z)$ and
$\varphi\mathcal J^{(2)}=Q$, with $\varphi$ the displayed multiplier; (F2)
and (F4) are stated there for each component. Numbering in \cite{VZ} refers
to the preprint version.

\begin{lemma}\label{lem:exp}
Let $a=j+q-m$. Then
\begin{align}
c(J_z(h,j,k,l,m,q))&=\sum_{\lambda=0}^{a}\binom a\lambda z^{j+k+\lambda}\,
c(I_z(h,\,j,\,k+\lambda,\,l+a-\lambda,\,m+\lambda)),\label{eq:E1}\\
c(J_z(h,j,k,l,m,q))&=\sum_{\lambda=0}^{a}\binom a\lambda z^{j+k}(1-z)^{\lambda}\,
c(I_z(h,\,j,\,k,\,l+\lambda,\,m)).\label{eq:E3}
\end{align}
\end{lemma}

\begin{proof}
Expand $(1-y+zy)^a$ as $((1-y)+zy)^a$ and as $(z+(1-z)(1-y))^a$. Each term
has the denominator and the pole $y=x/(x-z)$ of the corresponding
$I_z$-integrand. The powers of $z$ are
$z^{k-l-q+\lambda}\cdot z^{(l+a-\lambda)+(m+\lambda)}=z^{j+k+\lambda}$ and
$z^{k-l-q+a-\lambda}\cdot z^{(l+\lambda)+m}=z^{j+k}$. The identities hold
for each component, hence for $c$.
\end{proof}

Identity \eqref{eq:E1} is used in the proof of \cite[Lemma~2.1]{VZ} and in
\cite[(3.11)--(3.12)]{Pr}.

\section{The divisor}\label{sec:div}

\subsection{Local bounds}
For the tuple in \eqref{eq:GD}, (F1) gives $H_1=323N$, $H_2=322N$,
$\alpha=132N$, $\beta=107N$. Since $(1-z)^{\beta}=\pm3^{107N}$ at $z=4$,
\begin{equation}\label{eq:native}
D_N\,c(G_N)\in\Z^2,\qquad G_N=c_1(G_N)-\theta\,c_2(G_N).
\end{equation}

Fix $N\ge1$ and a prime $p>3$ with $p^2>430N$. Put $x=N/p$ and
$F_a=\fl{ax}$, and let $\br{\cdot}$ be the Iverson bracket. For
$0\le n\le430N$, $v_p(n!)=\fl{n/p}$ and $v_p(d_n)=\br{n\ge p}$.

\begin{lemma}\label{lem:transfer}
Let $t$ be admissible with all entries of $t$ and $u(t)$ in $[0,430N]$, and
let $s\in O(t)$. Then, for $i=1,2$,
\[
v_p\bigl(c_i(I_4(t))\bigr)\ \ge\ \sum_{\nu=1}^{5}\Bigl(\fl{\frac{t_\nu}{p}}-\fl{\frac{s_\nu}{p}}\Bigr)
-\br{H(s)\ge p}-\br{K(s)\ge p}.
\]
\end{lemma}

\begin{proof}
By (F4), $c(I_4(t))=T(t)T(s)^{-1}c(I_4(s))$. The entries of $s$ and $u(s)$
are those of $t$ and $u(t)$ in some order. Apply (F3) to $s$, where $z=4$ and
$z-1=3$ are units at $p$.
\end{proof}

\begin{lemma}\label{lem:gains}
Suppose $p<322N$. For $b\in\{190,215\}$ put $e_b=2b-165$ and
\begin{align*}
C&=2F_{240}-F_{265}-F_{215},\\
U_b&=2F_{240}+2F_b-F_{e_b}-F_{323}-F_{322},\\
V_b&=2F_b+F_{157}-F_{e_b}-F_{322}.
\end{align*}
Then, for $i=1,2$,
\begin{equation}\label{eq:M}
v_p\bigl(D_N\,c_i(G_N)\bigr)\ \ge\ M:=\max\{0,\,C,\,U_{190},\,U_{215},\,V_{190},\,V_{215}\}.
\end{equation}
If moreover $p>265N$, then $v_p(D_N\,c_i(G_N))\ge1$.
\end{lemma}

\begin{proof}
Here $v_p(D_N)=2$. By (F2), $c(G_N)$ is unchanged if \eqref{eq:GD} is
replaced by $(240N,240N,215N,157N,190N,165N)$. So \eqref{eq:E1}, applied to
\eqref{eq:GD} for $b=190$ and to this tuple for $b=215$, gives
\[
c(G_N)=\sum_{\lambda=0}^{bN}\binom{bN}{\lambda}4^{(240+b)N+\lambda}\,c(I_4(t)),
\]
where $t=(240,240,b,157+b,405-b)N+(0,0,1,-1,1)\lambda$. Here
$u(t)=(322,\,645-2b,\,323-b,\,e_b)N+(0,0,1,0)\lambda$; all entries of $t$
and $u(t)$ lie in $[0,430N]$. The tuples
\begin{align*}
s_1=\sigma t&=(645-2b,\ e_b,\ 405-b,\ 157+b,\ b)N+(0,0,1,-1,1)\lambda,\\
s_2=\sigma\rho\sigma t&=(323-b,\ 322,\ e_b,\ 405-b,\ 157+b)N+(1,0,0,1,-1)\lambda,\\
s_3=\sigma\rho\sigma\rho\sigma t&=(240,\ 240,\ 322,\ e_b,\ 405-b)N+(0,0,0,0,1)\lambda
\end{align*}
lie in $O(t)$. Put $\tau=\lambda/p$. Then
$v_p\binom{bN}{\lambda}=F_b-\fl{\tau}-\fl{bx-\tau}$. Also
$\br{H(s)\ge p}+\br{K(s)\ge p}\le2=v_p(D_N)$ for each $s\in O(t)$, so
Lemma~\ref{lem:transfer} shows that the $\lambda$-th term of $D_Nc_i(G_N)$
has valuation at least
\[
g(s)=\sum_\nu\bigl(\fl{t_\nu/p}-\fl{s_\nu/p}\bigr)+F_b-\fl\tau-\fl{bx-\tau}.
\]
Direct subtraction gives
\begin{align*}
g(s_1)&=C+\bigl(F_b-\fl{\tau}-\fl{bx-\tau}\bigr),\\
g(s_2)&=2F_{240}+F_b-F_{e_b}-F_{322}+\bigl(\fl{bx+\tau}-\fl{\tau}\bigr)\\
&\qquad-\bigl(\fl{bx-\tau}+\fl{(323-b)x+\tau}\bigr),\\
g(s_3)&=F_b-F_{e_b}-F_{322}+\bigl(\fl{bx+\tau}-\fl{\tau}\bigr)\\
&\qquad+\bigl(\fl{(157+b)x-\tau}-\fl{bx-\tau}\bigr).
\end{align*}
By $\fl{y}+\fl{w}\le\fl{y+w}$,
\begin{gather*}
F_b-\fl\tau-\fl{bx-\tau}\ge0,\qquad \fl{bx+\tau}-\fl\tau\ge F_b,\\
\fl{(157+b)x-\tau}-\fl{bx-\tau}\ge F_{157},\qquad
\fl{bx-\tau}+\fl{(323-b)x+\tau}\le F_{323}.
\end{gather*}
Hence $g(s_1)\ge C$, $g(s_2)\ge U_b$ and $g(s_3)\ge V_b$ for every $\lambda$.
The ultrametric inequality and \eqref{eq:native} give \eqref{eq:M}.

Now let $p>265N$. By \eqref{eq:E3},
\[
c(G_N)=\sum_{\lambda=0}^{190N}\binom{190N}{\lambda}4^{430N}(-3)^{\lambda}\,c(I_4(t')),
\]
where
\begin{align*}
t'&=(240,240,190,157,215)N+(0,0,0,1,0)\lambda,\\
u(t')&=(132,265,323,215)N+(1,0,-1,0)\lambda,\\
s'=\sigma\rho\sigma t'&=(323,132,215,215,157)N+(-1,1,0,0,1)\lambda,\\
u(s')&=(240,265,240,190)N.
\end{align*}
So $\br{H(s')\ge p}=\br{K(s')\ge p}=0$. Also
$v_p\binom{190N}{\lambda}=0$ and $F_{240}=F_{215}=F_{190}=0$. By
Lemma~\ref{lem:transfer} and $v_p(D_N)=2$, each term of $D_Nc_i(G_N)$ has
valuation at least $2-\fl{323x-\tau}-\fl{132x+\tau}\ge2-\fl{455x}\ge1$, as
$455x<2$.
\end{proof}

\subsection{Periodic intervals}
Since $5$ divides $240$, $265$ and $215$, and $2\cdot240=265+215$, we get
$C(x+1/5)=C(x)$. For $1\le k\le47$ put $j_k=\fl{5k/48}$,
$\varrho_k=5k-48j_k$, $a_k=k/240$ and
\[
b_k=\begin{cases}(k-j_k)/215,&\varrho_k\le21,\\ (k+j_k+1)/265,&\varrho_k\ge22.\end{cases}
\]

\begin{lemma}\label{lem:band}
For $1\le k\le47$ one has $a_k<b_k\le(k+1)/240$ and $C=1$ on $(a_k,b_k)$.
\end{lemma}

\begin{proof}
Here $0<\varrho_k<48$, $(k-j_k)/215-a_k=\varrho_k/10320$,
$(k+j_k+1)/265-a_k=(48-\varrho_k)/12720$ and
$265(k-j_k)-215(k+j_k+1)=10\varrho_k-215$. So
$b_k=\min\{(k-j_k)/215,\,(k+j_k+1)/265\}>a_k$, and $48j_k\le5k$ gives
$b_k\le(k+j_k+1)/265\le(k+1)/240$. For $a_k<x<b_k$,
\[
k-j_k-1<k-\tfrac{5k}{48}<215x<k-j_k,\qquad
k+j_k\le k+\tfrac{5k}{48}<265x<k+j_k+1,
\]
and $F_{240}=k$. Hence $C=2k-(k-j_k-1)-(k+j_k)=1$.
\end{proof}

Put
$\mathcal C=\bigcup_{r=0}^{9}\bigcup_{k=1}^{47}\,(r/5+a_k,\ r/5+b_k)\subset(1/240,2)$;
by Lemma~\ref{lem:band} these $470$ intervals are pairwise disjoint. Let
$\mathcal E_1,\dots,\mathcal E_{18}$ be the open intervals
\[
\begin{array}{llll}
(1/322,\,1/265) & (1/215,\,2/323) & (2/215,\,3/322) & (3/215,\,4/265),\\
(4/190,\,7/323) & (7/265,\,9/322) & (5/157,\,11/323) & (7/190,\,12/322),\\
(9/215,\,14/323) & (11/215,\,17/322) & (11/190,\,19/323) & (17/265,\,21/322),\\
(11/157,\,23/323) & (14/190,\,24/322) & (15/190,\,26/323) & (19/215,\,29/322),\\
(21/190,\,36/322) & (22/190,\,38/323).
\end{array}
\]
They are pairwise disjoint and lie in $(1/322,1/5)$. Define
\[
\mu(x)=\mathbf 1_{\mathcal C}(x)+\sum_{i=1}^{18}\mathbf 1_{\mathcal E_i}(x).
\]

\begin{proposition}\label{prop:div}
Let $N\ge1$ and let $p>3$ be a prime with $p^2>430N$. Then
$\mu(N/p)\le v_p(D_N)$ and $v_p(D_N\,c_i(G_N))\ge\mu(N/p)$ for $i=1,2$.
\end{proposition}

\begin{proof}
Let $x=N/p$ with $\mu(x)>0$. Then $1/322<x<2$, so $p<322N$, $v_p(D_N)=2$
and $\mu(x)\le2$. If $x\notin\bigcup_i\mathcal E_i$, then
$\mu(x)=\mathbf 1_{\mathcal C}(x)\le\br{C=1}\le M$ by
Lemmas~\ref{lem:band} and~\ref{lem:gains}. If $x\in\mathcal E_1$, then
$x\notin\mathcal C$ and $p>265N$, so Lemma~\ref{lem:gains} gives
$v_p(D_N\,c_i(G_N))\ge1=\mu(x)$. If $x\in\mathcal E_2\cup\dots\cup\mathcal E_{18}$,
then $M\ge1+\br{C=1}\ge\mu(x)$ by Table~\ref{tab:floors}, which lists the
floors $F_a$ for $a\in\{157,190,215,240,265,322,323\}$ on pieces with no
jump inside.
\end{proof}

\section{Density}\label{sec:dens}

\begin{lemma}\label{lem:dens}
Let $\Delta_N=\prod p^{\mu(N/p)}$ over primes $p>3$ with $p^2>430N$. Then
$\Delta_N\mid D_N$, $(D_N/\Delta_N)\,c(G_N)\in\Z^2$, and
\begin{align*}
\lim_{N\to\infty}\frac{\log\Delta_N}{40N}&=\delta:=\frac1{40}\int_0^2\mu(x)\,\frac{dx}{x^2}>4.7670,\\
\lim_{N\to\infty}\frac{\log D_N}{40N}&=c_D:=\frac{645+132\log4+107\log3}{40}<23.63856.
\end{align*}
\end{lemma}

\begin{proof}
The first two claims follow from Proposition~\ref{prop:div} and
\eqref{eq:native}. If $\mu(N/p)>0$, then $p>N/2$, so $p>3$ and $p^2>430N$
for $N>1720$. For each of the $488$ intervals $I$ of $\mathcal C$
and $\mathcal E_1,\dots,\mathcal E_{18}$, the prime number theorem gives
\[
\sum_{N/p\in I}\log p=N\int_Ix^{-2}\,dx+o(N).
\]
Summing over $I$ gives $\delta$.
Likewise $\log d_n=n+o(n)$ gives $c_D$, and $c_D<23.63856$ by
Table~\ref{tab:num}.

The intervals of $\mathcal C$ with $r=0$ contribute
\[
\delta_0=\frac1{40}\sum_{k=1}^{47}\Bigl(\frac1{a_k}-\frac1{b_k}\Bigr)
=\frac1{40}\Bigl(\sum_{k=1}^{47}\frac{240}k-\sum_{\varrho_k\le21}\frac{215}{k-j_k}
-\sum_{\varrho_k\ge22}\frac{265}{k+j_k+1}\Bigr)>1.6187,
\]
with the two cases listed in Table~\ref{tab:band}. Let
$\ell=\sum_k(b_k-a_k)=1139/22790$ and
$\bar y=\ell^{-1}\sum_k\int_{a_k}^{b_k}y\,dy=164041/1640160$. For $r\ge1$,
Jensen's inequality for $y\mapsto(r/5+y)^{-2}$ gives
\[
\sum_k\Bigl(\frac1{r/5+a_k}-\frac1{r/5+b_k}\Bigr)
=\sum_k\int_{a_k}^{b_k}\frac{dy}{(r/5+y)^2}\ge\frac{\ell}{(r/5+\bar y)^2}.
\]
So the intervals with $1\le r\le9$ contribute at least
$\frac{\ell}{40}\sum_r(r/5+\bar y)^{-2}>0.0260$ (Table~\ref{tab:num}). The intervals $\mathcal E_i$
contribute $\frac{202007398913}{64696932300}>3.1223$. Thus
$\delta>1.6187+0.0260+3.1223=4.7670$.
\end{proof}

\section{Analytic estimates}\label{sec:an}

\subsection{The double integral}
Write $J_N^{(0)}$, $J_N^{(1)}$ for $J_4^{(0)}$, $J_4^{(1)}$ at the tuple
\eqref{eq:GD}.

\begin{lemma}\label{lem:subst}
Let $y=\eta/(\eta-4)$. Then $R_4(x,y)\,dy=-f_N(x,\eta)\,d\eta$, where
\[
f_N(x,\eta)=\frac{\Psi(x,\eta)^N}{(x-\eta)(4-\eta)},\qquad
\Psi(x,\eta)=\frac{x^{240}(1-x)^{240}\,\eta^{190}(1-\eta)^{190}}
{(x-\eta)^{215}(4-\eta)^{322}} .
\]
In particular, with
\[
\Phi(x,\tau)=\Psi(x,-\tau)=\frac{x^{240}(1-x)^{240}\,\tau^{190}(1+\tau)^{190}}
{(x+\tau)^{215}(4+\tau)^{322}},
\]
we have
\[
J_N^{(0)}=\int_0^1\!\!\int_0^\infty\frac{\Phi(x,\tau)^N}{(x+\tau)(4+\tau)}\,d\tau\,dx .
\]
\end{lemma}

\begin{proof}
We have
\begin{gather*}
y=\frac{-\eta}{4-\eta},\qquad 1-y=\frac4{4-\eta},\qquad
1+3y=\frac{4(1-\eta)}{4-\eta},\\
x(1-y)+4y=\frac{4(x-\eta)}{4-\eta},\qquad dy=\frac{-4\,d\eta}{(4-\eta)^2}.
\end{gather*}
The exponents are $j=h=240N$, $k=190N$, $l=157N$, $j+q-m=190N$,
$j+k-m+1=215N+1$ and $k-l-q=-132N$. The powers of $4$ cancel, since
$-132+157+190-215=0$ and $4^{-1}\cdot4=1$. Those of $4-\eta$ add to
$-(322N+1)$, and $(-\eta)^{190N}=\eta^{190N}$. For $J_N^{(0)}$ put
$\tau=-\eta$. Then $\tau$ runs over $[0,\infty)$ as $y$ runs over $[0,1)$,
and $-f_N\,d\eta=f_N(x,-\tau)\,d\tau$.
\end{proof}

\begin{lemma}\label{lem:real}
For $N\ge1$, $10^{-9}e^{-40(18.88)N}<J_N^{(0)}<e^{-40(18.87333)N}$.
\end{lemma}

\begin{proof}
Weighted AM--GM, $x+\tau\ge215\,(x/25)^{25/215}(\tau/190)^{190/215}$, gives
\begin{align*}
\Phi&\le\Bigl(\frac5{43}\Bigr)^{25}\Bigl(\frac{38}{43}\Bigr)^{190}
x^{215}(1-x)^{240}\,\frac{(1+\tau)^{190}}{(4+\tau)^{322}}\\
&\le\Phi_+:=\Bigl(\frac5{43}\Bigr)^{25}\Bigl(\frac{38}{43}\Bigr)^{190}
\Bigl(\frac{43}{91}\Bigr)^{215}\Bigl(\frac{48}{91}\Bigr)^{240}
\frac{(95/22)^{190}}{(161/22)^{322}},
\end{align*}
with maxima at $x=43/91$ and $\tau=73/22$. By Table~\ref{tab:num},
$-\log\Phi_+>40(18.87333)$. Also
\[
\int_0^1\!\!\int_0^\infty\frac{d\tau\,dx}{(x+\tau)(4+\tau)}=\int_0^1\frac{\log(4/x)}{4-x}\,dx
<\frac{1+\log4}{3}<1,
\]
so $J_N^{(0)}<\Phi_+^N$.
For the lower bound let $\mathcal B=[0.4711,0.4713]\times[3.3999,3.4001]$.
Bounding each factor of $\Phi$ by its least value on $\mathcal B$ gives
\[
\Phi\ge\Phi_-:=\frac{0.4711^{240}\,0.5287^{240}\,3.3999^{190}\,4.3999^{190}}
{3.8714^{215}\,7.4001^{322}}\qquad\text{on }\mathcal B,
\]
and $(x+\tau)(4+\tau)<29$ there. By Table~\ref{tab:num},
$-\log\Phi_-<40(18.88)$. Hence
$J_N^{(0)}\ge\operatorname{area}(\mathcal B)\,\Phi_-^N/29=4\cdot10^{-8}\,\Phi_-^N/29>10^{-9}e^{-40(18.88)N}$.
\end{proof}

\subsection{The residue integral}
Put $A=(3+4i)/5$, $B=(5-12i)/13$ and, for $0\le t,u\le1$,
\[
X(t)=\frac{t}{t+A(1-t)},\qquad Y(u)=\frac{u}{u+B(1-u)},\qquad
\omega=t(1-u)+u(1-t).
\]

\begin{lemma}\label{lem:arcs}
For $0\le t,u\le1$ and $x\in[0,1]$:
\begin{enumerate}
\item $|X'|\le5/4$, $|Y'|\le13/9$, $|X(1-X)|\le\frac54t(1-t)\le\frac54\omega$,
$|Y(1-Y)|\le\frac{13}9u(1-u)\le\frac{13}9\omega$, $|Y|\le1$ and $|4-Y|\ge3$;
\item $|X-Y|\ge7\omega/\sqrt{65}$, and $|x-Y|\ge2\omega_x/\sqrt{13}$ with
$\omega_x=x(1-u)+u(1-x)$;
\item $\operatorname{Im}X(t)<0<\operatorname{Im}Y(u)$ for $0<t,u<1$.
\end{enumerate}
\end{lemma}

\begin{proof}
One has $|t+A(1-t)|^2=1-\frac45t(1-t)\ge\frac45$,
$|u+B(1-u)|^2=1-\frac{16}{13}u(1-u)\ge\frac9{13}$,
$X'=A/(t+A(1-t))^2$ and $1-X=A(1-t)/(t+A(1-t))$; similarly for $Y$.
Also $|u+B(1-u)|^2-u^2=(1-u)(13-3u)/13\ge0$, so $|Y|\le1$. Next
$t(1-t)\le\omega$, since $\omega-t(1-t)$ equals $u(1-2t)+t^2\ge0$ for
$t\le\frac12$ and is at least $(1-t)^2$ for $t\ge\frac12$; likewise
$u(1-u)\le\omega$. For (2),
$X-Y=\frac{tB(1-u)-uA(1-t)}{(t+A(1-t))(u+B(1-u))}$ and the denominator has
modulus at most $1$. With $\xi=t(1-u)$, $\zeta=u(1-t)$ and
$\gamma=\operatorname{Re}(B\bar A)=-33/65$,
\[
|tB(1-u)-uA(1-t)|^2=\xi^2+\zeta^2-2\gamma\xi\zeta
=\tfrac{1-\gamma}2\,\omega^2+\tfrac{1+\gamma}2(\xi-\zeta)^2\ge\tfrac{49}{65}\,\omega^2 .
\]
For real $x$ take $A=1$, $t=x$ and $\gamma=\operatorname{Re}B=5/13$. Part (3)
follows from $\operatorname{Im}X=-\frac45t(1-t)/|t+A(1-t)|^2$ and its
analogue.
\end{proof}

\begin{lemma}\label{lem:contour}
For $N\ge1$, $J_N^{(1)}=\pi^{-1}\operatorname{Im}Z_N$, where
\[
Z_N=\int_0^1\!\!\int_0^1f_N\bigl(X(t),Y(u)\bigr)\,X'(t)\,Y'(u)\,dt\,du .
\]
\end{lemma}

\begin{proof}
Residues are invariant under the change of variables of
Lemma~\ref{lem:subst}, and $y=x/(x-4)$ corresponds to $\eta=x$. Hence
$J_N^{(1)}=-\int_0^1\Res_{\eta=x}f_N(x,\eta)\,dx$. Fix $x\in(0,1)$. By
Lemma~\ref{lem:arcs}, $Y$ runs from $0$ to $1$ in the upper half-plane and
$\bar Y$ in the lower one, both in the closed unit disc. So $\bar Y$ followed
by $Y$ reversed is a positively oriented Jordan curve around $x$, and the
other pole $\eta=4$ lies outside. As $f_N$ has real coefficients, the residue
theorem gives
\[
\Res_{\eta=x}f_N=\frac{\overline{W(x)}-W(x)}{2\pi i}=-\frac{\operatorname{Im}W(x)}{\pi},
\qquad W(x)=\int_0^1f_N(x,Y(u))\,Y'(u)\,du .
\]
By Lemma~\ref{lem:arcs}, also with $t$ replaced by $x$,
\[
|f_N(x,Y)\,Y'|\ll_N\omega_x^{215N-1},\qquad |f_N(X,Y)\,X'Y'|\ll_N\omega^{215N-1},
\]
so both are bounded on $[0,1]^2$. For
fixed $u\in(0,1)$ the function $w\mapsto f_N(w,Y(u))$ is holomorphic off
$w=Y(u)$, and $\operatorname{Im}Y(u)>0\ge\operatorname{Im}X(t)$. Cauchy's
theorem moves $\int_0^1dx$ to the arc $X$, and Fubini's theorem gives
$\int_0^1W(x)\,dx=Z_N$.
\end{proof}

\begin{lemma}\label{lem:res}
For $N\ge1$, $|J_N^{(1)}|<e^{-40(18.8995)N}$.
\end{lemma}

\begin{proof}
Put $r=t/(1-t)=e^R$ and $s=u/(1-u)=e^S$, so $X=r/(r+A)$ and $Y=s/(s+B)$.
With
\[
Q_1=r^2+\tfrac65r+1,\quad Q_2=s^2+\tfrac{10}{13}s+1,\quad
Q_3=r^2+\tfrac{66}{65}rs+s^2,\quad Q_4=9s^2+\tfrac{120}{13}s+16,
\]
one has $|X(1-X)|=r/Q_1$, $|Y(1-Y)|=s/Q_2$, $|X-Y|^2=Q_3/(Q_1Q_2)$ and
$|4-Y|^2=Q_4/Q_2$. So $\log|\Psi(X,Y)|=L$, where
\[
L=240\log r+190\log s-\tfrac{265}2\log Q_1+\tfrac{157}2\log Q_2
-\tfrac{215}2\log Q_3-161\log Q_4 .
\]
Let $\ell_\gamma(v)=\log(\cosh v+\gamma)$ and $\kappa=\log(4/3)$. Since
$Q_1=2r(\cosh R+\frac35)$, $Q_2=2s(\cosh S+\frac5{13})$,
$Q_3=2rs(\cosh(R-S)+\frac{33}{65})$ and $Q_4=24s(\cosh(S-\kappa)+\frac5{13})$,
\begin{align*}
L={}&-\tfrac{323}2\log2-161\log24-\tfrac{265}2\ell_{3/5}(R)+\tfrac{157}2\ell_{5/13}(S)\\
&-\tfrac{215}2\ell_{33/65}(R-S)-161\,\ell_{5/13}(S-\kappa).
\end{align*}
Here $\ell_\gamma''=(1+\gamma\cosh v)/(\cosh v+\gamma)^2>0$ and, for
$0<\gamma<1$,
\[
(\log\ell_\gamma'')'=-\sinh v\,\frac{2+\gamma\cosh v-\gamma^2}{(1+\gamma\cosh v)(\cosh v+\gamma)},
\qquad
|(\log\ell_\gamma'')'|\le\frac{2|\sinh v|}{\cosh v+\gamma}<2,
\]
so $\ell_\gamma''(v-\kappa)\ge\frac9{16}\ell_\gamma''(v)$. Therefore
\begin{align*}
-d^2L={}&\tfrac{265}2\ell_{3/5}''(R)\,dR^2+\tfrac{215}2\ell_{33/65}''(R-S)(dR-dS)^2\\
&+\bigl(161\,\ell_{5/13}''(S-\kappa)-\tfrac{157}2\ell_{5/13}''(S)\bigr)dS^2
\end{align*}
is positive definite, since $161\cdot\frac9{16}>\frac{157}2$. So $L$ is
strictly concave and $L_{RS}>0$. The partial derivatives are
\begin{align*}
L_R&=240-265\,\frac{r(r+3/5)}{Q_1}-215\,\frac{r(r+33s/65)}{Q_3},\\
L_S&=190+157\,\frac{s(s+5/13)}{Q_2}-215\,\frac{s(s+33r/65)}{Q_3}-322\,\frac{s(9s+60/13)}{Q_4} .
\end{align*}
Table~\ref{tab:grad} gives $L_R,L_S>0$ at $(r,s)=(\frac{29}{25},\frac{139}{100})$
and $L_R,L_S<0$ at $(\frac{59}{50},\frac{71}{50})$. Consider the box
$\frac{29}{25}\le r\le\frac{59}{50}$, $\frac{139}{100}\le s\le\frac{71}{50}$.
Since $L_{RS}>0$, we get $L_R>0$ on its face $r=\frac{29}{25}$, $L_R<0$ on
$r=\frac{59}{50}$, $L_S>0$ on $s=\frac{139}{100}$ and $L_S<0$ on
$s=\frac{71}{50}$.
So the maximum of $L$ on the box is an interior critical point, hence the
global maximum. At $(r,s)=(\frac{117}{100},\frac75)$,
Table~\ref{tab:grad} gives $-\frac14<L_R<0<L_S<\frac3{10}$ and $L<-755.99$.
Concavity gives
\[
\sup L<-755.99+\tfrac14\log\tfrac{117}{116}+\tfrac3{10}\log\tfrac{71}{70}
<-755.99+\tfrac1{464}+\tfrac3{700}<-755.98 .
\]
By Lemma~\ref{lem:arcs}, $|X'Y'|<2$, $|X-Y|\ge\frac45\omega$ and $|4-Y|\ge3$.
With $w=1-2u$,
\[
\int_0^1\!\!\int_0^1\frac{dt\,du}{\omega}=\int_0^1\frac{\log\frac{1-u}{u}}{1-2u}\,du
=\frac12\int_{-1}^1\log\frac{1+w}{1-w}\,\frac{dw}{w}
=2\sum_{n\ge0}\frac1{(2n+1)^2}=\frac{\pi^2}4 .
\]
By Lemma~\ref{lem:contour},
\[
|J_N^{(1)}|\le\frac1\pi\cdot2\cdot\frac54\cdot\frac13\cdot\frac{\pi^2}4\,e^{-40(18.8995)N}
=\frac{5\pi}{24}\,e^{-40(18.8995)N}<e^{-40(18.8995)N}.\qedhere
\]
\end{proof}

\section{Proof of Theorem~\ref{thm:main}}\label{sec:proof}

By Lemmas~\ref{lem:real} and~\ref{lem:res},
$J_N^{(0)}>10^{-9}e^{-40(18.88)N}>(\log4)\,|J_N^{(1)}|$ for large $N$.
Hence $G_N=J_N^{(0)}-(\log4)J_N^{(1)}>0$ and $G_N<3e^{-40(18.87333)N}$. Put
\[
\Lambda_N=\frac{D_N}{\Delta_N}\,G_N
=\frac{D_N}{\Delta_N}\,c_1(G_N)-\frac{D_N}{\Delta_N}\,c_2(G_N)\,\theta .
\]
By \eqref{eq:native} and Lemma~\ref{lem:dens}, $\Lambda_N\in\Z+\Z\theta$.
Also $\Lambda_N>0$ for large $N$, and
\[
\limsup_{N\to\infty}\frac{\log\Lambda_N}{40N}\le c_D-\delta-18.87333
<23.63856-4.7670-18.87333=-0.00177 .
\]
If $q_0\theta\in\Z$ for an integer $q_0\ge1$, then $q_0\Lambda_N$ is a
positive integer for large $N$ and $q_0\Lambda_N\to0$, which is impossible.
\qed

\appendix

\section{Floor data}\label{app:floor}

Table~\ref{tab:floors} cuts each $\mathcal E_i$ at the jumps of the $F_a$.
On each piece, $F_a=\fl{ax_0}$ with $x_0$ the left end. The columns $C$,
$U_b$, $V_b$ follow from Lemma~\ref{lem:gains}, and the last column is
$1+\br{C=1}$. The first row uses \eqref{eq:E3} instead of $C$, $U_b$, $V_b$.

{\footnotesize
\setlength{\tabcolsep}{2.2pt}
\begin{longtable}{l|*{7}{>{$}r<{$}}|*{5}{>{$}r<{$}}|>{$}c<{$}}
\caption{Floors on the pieces of $\mathcal E_1,\dots,\mathcal E_{18}$.}\label{tab:floors}\\
\text{piece} & F_{157} & F_{190} & F_{215} & F_{240} & F_{265} & F_{322} & F_{323}
& C & U_{190} & U_{215} & V_{190} & V_{215} & 1{+}\br{C{=}1}\\ \hline
\endfirsthead
\text{piece} & F_{157} & F_{190} & F_{215} & F_{240} & F_{265} & F_{322} & F_{323}
& C & U_{190} & U_{215} & V_{190} & V_{215} & 1{+}\br{C{=}1}\\ \hline
\endhead
$(1/322,1/265)$ & 0 & 0 & 0 & 0 & 0 & 1 & 1 & \multicolumn{5}{c|}{by \eqref{eq:E3}} & 1 \\
\hline
$(1/215,1/190)$ & 0 & 0 & 1 & 1 & 1 & 1 & 1 & 0 & -1 & 1 & -2 & 0 & 1 \\
$[1/190,2/323)$ & 0 & 1 & 1 & 1 & 1 & 1 & 1 & 0 & 1 & 1 & 0 & 0 & 1 \\
\hline
$(2/215,3/322)$ & 1 & 1 & 2 & 2 & 2 & 2 & 3 & 0 & -1 & 1 & -1 & 1 & 1 \\
\hline
$(3/215,4/265)$ & 2 & 2 & 3 & 3 & 3 & 4 & 4 & 0 & -1 & 1 & -1 & 1 & 1 \\
\hline
$(4/190,7/323)$ & 3 & 4 & 4 & 5 & 5 & 6 & 6 & 1 & 2 & 1 & 1 & 0 & 2 \\
\hline
$(7/265,9/323)$ & 4 & 5 & 5 & 6 & 7 & 8 & 8 & 0 & 1 & -1 & 1 & -1 & 1 \\
$[9/323,6/215)$ & 4 & 5 & 5 & 6 & 7 & 8 & 9 & 0 & 0 & -2 & 1 & -1 & 1 \\
$[6/215,9/322)$ & 4 & 5 & 6 & 6 & 7 & 8 & 9 & -1 & -1 & 0 & 0 & 1 & 1 \\
\hline
$(5/157,7/215)$ & 5 & 6 & 6 & 7 & 8 & 10 & 10 & 0 & 0 & -2 & 1 & -1 & 1 \\
$[7/215,8/240)$ & 5 & 6 & 7 & 7 & 8 & 10 & 10 & -1 & -1 & 0 & 0 & 1 & 1 \\
$[8/240,9/265)$ & 5 & 6 & 7 & 8 & 8 & 10 & 10 & 1 & 1 & 2 & 0 & 1 & 2 \\
$[9/265,11/323)$ & 5 & 6 & 7 & 8 & 9 & 10 & 10 & 0 & 1 & 1 & 0 & 0 & 1 \\
\hline
$(7/190,12/323)$ & 5 & 7 & 7 & 8 & 9 & 11 & 11 & 0 & 1 & -1 & 1 & -1 & 1 \\
$[12/323,8/215)$ & 5 & 7 & 7 & 8 & 9 & 11 & 12 & 0 & 0 & -2 & 1 & -1 & 1 \\
$[8/215,12/322)$ & 5 & 7 & 8 & 8 & 9 & 11 & 12 & -1 & -1 & 0 & 0 & 1 & 1 \\
\hline
$(9/215,8/190)$ & 6 & 7 & 9 & 10 & 11 & 13 & 13 & 0 & -1 & 1 & -2 & 0 & 1 \\
$[8/190,14/323)$ & 6 & 8 & 9 & 10 & 11 & 13 & 13 & 0 & 1 & 1 & 0 & 0 & 1 \\
\hline
$(11/215,10/190)$ & 8 & 9 & 11 & 12 & 13 & 16 & 16 & 0 & -1 & 1 & -1 & 1 & 1 \\
$[10/190,17/322)$ & 8 & 10 & 11 & 12 & 13 & 16 & 17 & 0 & 0 & 0 & 1 & 1 & 1 \\
\hline
$(11/190,14/240)$ & 9 & 11 & 12 & 13 & 15 & 18 & 18 & -1 & 0 & -1 & 1 & 0 & 1 \\
$[14/240,19/323)$ & 9 & 11 & 12 & 14 & 15 & 18 & 18 & 1 & 2 & 1 & 1 & 0 & 2 \\
\hline
$(17/265,21/323)$ & 10 & 12 & 13 & 15 & 17 & 20 & 20 & 0 & 1 & -1 & 1 & -1 & 1 \\
$[21/323,14/215)$ & 10 & 12 & 13 & 15 & 17 & 20 & 21 & 0 & 0 & -2 & 1 & -1 & 1 \\
$[14/215,21/322)$ & 10 & 12 & 14 & 15 & 17 & 20 & 21 & -1 & -1 & 0 & 0 & 1 & 1 \\
\hline
$(11/157,17/240)$ & 11 & 13 & 15 & 16 & 18 & 22 & 22 & -1 & -1 & 0 & 0 & 1 & 1 \\
$[17/240,23/323)$ & 11 & 13 & 15 & 17 & 18 & 22 & 22 & 1 & 1 & 2 & 0 & 1 & 2 \\
\hline
$(14/190,24/323)$ & 11 & 14 & 15 & 17 & 19 & 23 & 23 & 0 & 1 & -1 & 1 & -1 & 1 \\
$[24/323,16/215)$ & 11 & 14 & 15 & 17 & 19 & 23 & 24 & 0 & 0 & -2 & 1 & -1 & 1 \\
$[16/215,24/322)$ & 11 & 14 & 16 & 17 & 19 & 23 & 24 & -1 & -1 & 0 & 0 & 1 & 1 \\
\hline
$(15/190,17/215)$ & 12 & 15 & 16 & 18 & 20 & 25 & 25 & 0 & 0 & -2 & 1 & -1 & 1 \\
$[17/215,19/240)$ & 12 & 15 & 17 & 18 & 20 & 25 & 25 & -1 & -1 & 0 & 0 & 1 & 1 \\
$[19/240,21/265)$ & 12 & 15 & 17 & 19 & 20 & 25 & 25 & 1 & 1 & 2 & 0 & 1 & 2 \\
$[21/265,26/323)$ & 12 & 15 & 17 & 19 & 21 & 25 & 25 & 0 & 1 & 1 & 0 & 0 & 1 \\
\hline
$(19/215,14/157)$ & 13 & 16 & 19 & 21 & 23 & 28 & 28 & 0 & -1 & 1 & -2 & 0 & 1 \\
$[14/157,17/190)$ & 14 & 16 & 19 & 21 & 23 & 28 & 28 & 0 & -1 & 1 & -1 & 1 & 1 \\
$[17/190,29/323)$ & 14 & 17 & 19 & 21 & 23 & 28 & 28 & 0 & 1 & 1 & 1 & 1 & 1 \\
$[29/323,29/322)$ & 14 & 17 & 19 & 21 & 23 & 28 & 29 & 0 & 0 & 0 & 1 & 1 & 1 \\
\hline
$(21/190,36/323)$ & 17 & 21 & 23 & 26 & 29 & 35 & 35 & 0 & 1 & -1 & 1 & -1 & 1 \\
$[36/323,24/215)$ & 17 & 21 & 23 & 26 & 29 & 35 & 36 & 0 & 0 & -2 & 1 & -1 & 1 \\
$[24/215,36/322)$ & 17 & 21 & 24 & 26 & 29 & 35 & 36 & -1 & -1 & 0 & 0 & 1 & 1 \\
\hline
$(22/190,25/215)$ & 18 & 22 & 24 & 27 & 30 & 37 & 37 & 0 & 0 & -2 & 1 & -1 & 1 \\
$[25/215,28/240)$ & 18 & 22 & 25 & 27 & 30 & 37 & 37 & -1 & -1 & 0 & 0 & 1 & 1 \\
$[28/240,31/265)$ & 18 & 22 & 25 & 28 & 30 & 37 & 37 & 1 & 1 & 2 & 0 & 1 & 2 \\
$[31/265,38/323)$ & 18 & 22 & 25 & 28 & 31 & 37 & 37 & 0 & 1 & 1 & 0 & 0 & 1 \\
\end{longtable}}

\begin{table}[ht]
\caption{The two cases in the definition of $b_k$.}\label{tab:band}
\centering\small
\begin{tabular}{c|cc|cc}
$j_k$ & $\varrho_k\le21$ & $k-j_k$ & $\varrho_k\ge22$ & $k+j_k+1$\\ \hline
0 & $1\le k\le4$ & $1,\dots,4$ & $5\le k\le9$ & $6,\dots,10$\\
1 & $10\le k\le13$ & $9,\dots,12$ & $14\le k\le19$ & $16,\dots,21$\\
2 & $20\le k\le23$ & $18,\dots,21$ & $24\le k\le28$ & $27,\dots,31$\\
3 & $29\le k\le33$ & $26,\dots,30$ & $34\le k\le38$ & $38,\dots,42$\\
4 & $39\le k\le42$ & $35,\dots,38$ & $43\le k\le47$ & $48,\dots,52$
\end{tabular}
\end{table}

\section{Numerical values}\label{app:num}

Values are truncated to the digits shown.

\begin{table}[ht]
\caption{Constants.}\label{tab:num}
\centering\small
\begin{tabular}{l|l}
quantity & value\\ \hline
$-\log\Phi_+/40$ & $18.8733376741\ldots$\\
$-\log\Phi_-/40$ & $18.8775167011\ldots$\\
$c_D$ & $23.6385592638\ldots$\\
$\delta_0$ & $1.6187800851\ldots$\\
$\frac{\ell}{40}\sum_{r=1}^9(r/5+\bar y)^{-2}$ & $0.0260767666\ldots$\\
$\frac1{40}\sum_i\int_{\mathcal E_i}x^{-2}\,dx$ & $3.1223644109\ldots$
\end{tabular}
\end{table}

\begin{table}[ht]
\caption{Values of $L_R$, $L_S$ and $L$.}\label{tab:grad}
\centering\small
\begin{tabular}{c|ccc}
$(r,s)$ & $L_R$ & $L_S$ & $L$\\ \hline
$(29/25,\,139/100)$ & $0.57635\ldots$ & $0.62026\ldots$ & \\
$(59/50,\,71/50)$ & $-0.52544\ldots$ & $-1.02457\ldots$ & \\
$(117/100,\,7/5)$ & $-0.22659\ldots$ & $0.26833\ldots$ & $-755.99795\ldots$
\end{tabular}
\end{table}

\section{Computer checks}\label{app:ver}

The two scripts in Table~\ref{tab:checks} recheck only the statements listed
there, in exact rational arithmetic. They are included as ancillary files.
The proof does not use them. For
$\log v$ the analytic script writes $v=2^eq$ with $e\in\Z$, $1\le q\le2$, and
uses $30$ terms of $\log q=2\sum_{n\ge0}w^{2n+1}/(2n+1)$, $w=(q-1)/(q+1)$,
with tail at most $2w^{61}/(61(1-w^2))$.

\begin{table}[ht]
\caption{Scripts.}\label{tab:checks}
\centering\small
\begin{tabular}{>{\raggedright}p{0.27\textwidth}>{\raggedright}p{0.53\textwidth}r}
script & statements checked & time\tabularnewline \hline
\texttt{check\_arith.py} & the intervals $\mathcal E_i$ and Tables~\ref{tab:floors}, \ref{tab:band} as printed; Lemma~\ref{lem:band}; the tuples $s_1,s_2,s_3,s'$ and the floor sums in Lemma~\ref{lem:gains}; Proposition~\ref{prop:div} for $1/265\le x<2$ at each jump of the $F_a$ and between consecutive jumps, and $M=0$, $\mu=1$ on $\mathcal E_1$; the numbers in Lemma~\ref{lem:dens}; $57/40$ and $3.35$ in Section~1; rows 4--6 of Table~\ref{tab:num} & 10\,s\tabularnewline
\texttt{check\_arith.py --brute 1,2,3,4,5,6,7,8} & for $N\le8$ and all primes $p$ with $3<p\le323N$ and $p^2>430N$ (for larger $p$, $\mu(N/p)=0$), the best bound from Lemma~\ref{lem:transfer} over \eqref{eq:E1} with $b=190,215$ and \eqref{eq:E3}, all $\lambda$ and all $s\in O(t)$, is at least $\mu(N/p)$ & 3\,min\tabularnewline
\texttt{check\_analytic.py} & the exponents in Lemma~\ref{lem:subst}; the constants in Lemma~\ref{lem:real}; the identities in Lemmas~\ref{lem:arcs} and~\ref{lem:res} at rational points; rows 1--3 of Table~\ref{tab:num} and Table~\ref{tab:grad} as printed; $23.63856-4.7670-18.87333=-0.00177$ & 1\,s\tabularnewline
\end{tabular}
\end{table}

\section*{Use of AI tools}

Codex (OpenAI, GPT-5.6 Sol) searched for the parameters in \eqref{eq:GD}
and wrote the Python code and the interval lists. Claude Code (Anthropic,
Claude Fable 5) checked the code, the parameters and the tables. The author
checked every step of the proof, every table and every reference, and takes
full responsibility for the content of the paper.

\begin{thebibliography}{10}

\bibitem{CDT}
F.~Calegari, V.~Dimitrov, and Y.~Tang, \emph{The linear independence of $1$,
$\zeta(2)$, and $L(2,\chi_{-3})$}, preprint (2024), arXiv:2408.15403v2.

\bibitem{DLMF}
\emph{NIST Digital Library of Mathematical Functions},
\url{https://dlmf.nist.gov/}, Release 1.2.8 of 2026-09-15, F.~W.~J. Olver
et al., eds.

\bibitem{Ha}
M.~Hata, \emph{Rational approximations to the dilogarithm}, Trans. Amer. Math.
Soc. \textbf{336} (1993), no.~1, 363--387. MR1147401

\bibitem{Pr}
T.~Prellberg, \emph{Three irrationality results for the dilogarithm:
Rhin--Viola and Viola--Zudilin constructions at $-1/4$, $1/5$, and $-1/3$},
preprint (2026), arXiv:2609.06895v1.

\bibitem{RV96}
G.~Rhin and C.~Viola, \emph{On a permutation group related to $\zeta(2)$},
Acta Arith. \textbf{77} (1996), no.~1, 23--56. MR1404975

\bibitem{RV05}
G.~Rhin and C.~Viola, \emph{The permutation group method for the dilogarithm},
Ann. Sc. Norm. Super. Pisa Cl. Sci. (5) \textbf{4} (2005), no.~3, 389--437.
MR2185863

\bibitem{RV19}
G.~Rhin and C.~Viola, \emph{Linear independence of $1$, $\mathrm{Li}_1$ and
$\mathrm{Li}_2$}, Mosc. J. Comb. Number Theory \textbf{8} (2019), no.~1,
81--96. MR3864310

\bibitem{VZ}
C.~Viola and W.~Zudilin, \emph{Linear independence of dilogarithmic values},
J. Reine Angew. Math. \textbf{736} (2018), 193--223; preprint version:
Max-Planck-Institut f\"ur Mathematik Preprint Series 2014 (26),
\url{http://www.mpim-bonn.mpg.de/preblob/5515}. MR3769989

\end{thebibliography}

\end{document}
